% chiral_moonshine_unified.tex
%
% Chiral moonshine through signed denominator/bar comparison:
% denominator products compare with bar data only after the product
% exponents have been identified, with sign and normalization fixed,
% with the primitive graded supertrace of the chosen chiral object.

\providecommand{\cA}{\mathcal A}
\providecommand{\cB}{\mathcal B}
\providecommand{\cH}{\mathcal H}
\providecommand{\cM}{\mathcal M}
\providecommand{\cO}{\mathcal O}
\providecommand{\PhiFA}{\Phi^{\mathrm{FA}}}
\providecommand{\SpCh}{\mathrm{Sp}^{\mathrm{ch}}}
\providecommand{\cT}{\mathcal T}
\providecommand{\bZ}{\mathbb Z}
\providecommand{\bQ}{\mathbb Q}
\providecommand{\bH}{\mathbb H}
\providecommand{\bM}{\mathbb M}
\providecommand{\Aut}{\operatorname{Aut}}
\providecommand{\Tr}{\operatorname{Tr}}
\providecommand{\Str}{\operatorname{Str}}
\providecommand{\sdim}{\operatorname{sdim}}
\providecommand{\PE}{\operatorname{PE}}
\providecommand{\Co}{\mathrm{Co}}
\providecommand{\Th}{\mathrm{Th}}
\providecommand{\Mtw}{M_{24}}
\providecommand{\Vnat}{V^{\natural}}
\providecommand{\Vsnat}{V^{s\natural}}
\providecommand{\VTh}{V^{\Th}}
\providecommand{\VKthree}{V_{K3}^{\cN{=}4}}
\providecommand{\cN}{\mathcal N}
\providecommand{\Barch}{\mathrm B^{\mathrm{ch}}}
\providecommand{\BarchOrd}{\mathrm B^{\mathrm{ch},\mathrm{ord}}}
\providecommand{\kapch}{\kappa_{\mathrm{ch}}}
\providecommand{\kapBKM}{\kappa_{\mathrm{BKM}}}
\providecommand{\kapcat}{\kappa_{\mathrm{cat}}}
\providecommand{\kapfiber}{\kappa_{\mathrm{fiber}}}
\providecommand{\ZcSig}{Z^{\mathrm{der}}_{\mathrm{ch}}}
\providecommand{\ChirHoch}{\mathrm{ChirHoch}}

\chapter{Chiral moonshine: denominator products and primitive bar traces}
\label{chap:chiral-moonshine-unified}

\noindent
A Borcherds denominator is not a bar complex.  It is an automorphic
product whose exponents are root multiplicities of a generalized
Kac--Moody algebra.  A chiral bar complex supplies a product only
after its primitive weight strata have been converted into a
plethystic Euler product.  Moonshine enters when the same integers
occur on both sides.

The Monster scalar lane and the Monster-equivariant lane are different
objects.  The scalar lane of $\Vnat$ sees the Virasoro curvature
$\kapch(\Vnat)=12$, the Virasoro shadow coefficients at $c=24$, and
the collision residue
\[
 r^{\mathrm{Vir}}_{24}(z)=\frac{12}{z^3}+\frac{2T}{z}.
\]
The equivariant lane sees the Monster action on graded spaces and, for
a bar statement, on
$H^\bullet(\operatorname{Prim}\Barch(\Vnat))$.  The McKay--Thompson
series are not determined by the scalar tower alone.

Four sources carry the comparison.  The Monster source is the FLM
module $\Vnat$ and Borcherds' monster Lie algebra
\cite{FLM88,Bor92}.  The Conway source is Duncan's super-VOA
$\Vsnat$ with $\Co_0$ supertraces.  The Thompson source is the
Harvey--Rayhaun / Griffin--Mertens graded $\Th$-module, not an
unconditional VOA datum in the input package.  The Mathieu source is
the K3 elliptic genus, the Gannon $\Mtw$-module, and the
DMVV--Gritsenko--Nikulin multiplicative lift
\cite{Gannon16,DMVV,GritsenkoNikulin97}.

%%%%%%%%%%%%%%%%%%%%%%%%%%%%%%%%%%%%%%%%%%%%%%%%%%%%%%%%%%%%%%%%%%%%%%%%%%%%
\section{The four moonshine inputs}
\label{sec:moonshine-setup}
%%%%%%%%%%%%%%%%%%%%%%%%%%%%%%%%%%%%%%%%%%%%%%%%%%%%%%%%%%%%%%%%%%%%%%%%%%%%

\begin{definition}[Moonshine input package]
\label{def:moonshine-four-families}
The ordered package
\[
 \cM
 \;=\;
 \bigl(\Vnat,\;\Vsnat,\;W^{\Th},\;\VKthree\bigr)
\]
has the following meaning.
\begin{enumerate}[label=\textup{(\arabic*)}]
\item $\Vnat$ is the Frenkel--Lepowsky--Meurman holomorphic VOA of
central charge $24$, $\dim \Vnat_1=0$, with
$\Aut(\Vnat)=\bM$.
\item $\Vsnat$ is Duncan's $N{=}1$ Conway super-VOA of central charge
$12$, with $\dim \Vsnat_{1/2}=0$ and $\Co_0=2\cdot\Co_1$ acting by
super-VOA automorphisms.
\item $W^{\Th}$ denotes the graded Thompson supermodule whose graded
traces are the Harvey--Rayhaun Thompson moonshine functions.  The
notation $\VTh$ is reserved for a chiral algebra model realizing
$W^{\Th}$ if such a model is supplied.
\item $\VKthree$ denotes the $c=6$ $\cN{=}4$ chiral sector appearing
in the K3 elliptic genus.  Gannon's theorem supplies a graded
$\Mtw$-module for the Mathieu coefficients; it does not assert that
the full group $\Mtw$ acts as symmetries of one K3 sigma model.
\end{enumerate}
\end{definition}

\begin{remark}[Separated modular characteristics]
\label{rem:moonshine-kappa-ch-table}
The symbol $\kappa$ is lane-dependent:
\[
\begin{array}{c|c|c}
\text{object} & \text{invariant} & \text{value}\\
\hline
\Vnat & \kapch(\Vnat) & 12\\
V_{\Lambda_{24}}\text{ on the lattice lane}
 & \kappa_{\mathrm{latt}}(V_{\Lambda_{24}}) & 24\\
\Vsnat & \kapch(\Vsnat) & 6\\
\VKthree\text{ on the Virasoro lane} & \kapch(\VKthree) & 3\\
K3\text{ as a CY}_2\text{ category} & \kapcat(K3) & 2\\
\mathfrak g_{\Delta_5}\text{ BKM denominator} & \kapBKM(\mathfrak g_{\Delta_5}) & 5\\
\mathcal H_{\mathrm{Muk}}(K3)\text{ conductor lane} & K^\kappa & 8\\
\widetilde H(K3,\bZ)\text{ Mukai lattice} & \kapfiber & 24
\end{array}
\]
The Monster value is the proved scalar modular-characteristic statement
of Proposition~\ref{prop:moonshine-kappa}: the weight-$1$ Heisenberg
channel is absent and the Virasoro contribution gives $c/2=12$.  The
Leech value is the lattice formula
$\kappa_{\mathrm{latt}}(V_{\Lambda_{24}})=\operatorname{rank}(\Lambda_{24})=24$.
The Conway value is the same Virasoro calculation at $c=12$.  The K3
values live in different lanes.  The equality
$\kapch(\VKthree)=3$ is the $c/2$ Virasoro value for the $c=6$
$\cN{=}4$ algebra; it is not the categorical value
$\kapcat(K3)=\chi(\cO_{K3})=2$, not the Borcherds weight
$\kapBKM(\mathfrak g_{\Delta_5})=5$, not the scalar conductor
$K^\kappa=8$ of the Mukai-enhanced $\mathcal B$-family
\textup{(}Proposition~\ref{prop:archetype-complementarity-bridge};
$K^\kappa=2c_+(\mathrm{Mukai}(K3))$\textup{)}, and not the Mukai
rank $24$.
\end{remark}

\begin{remark}[Schellekens and Niemeier scope]
\label{rem:moonshine-schellekens-niemeier-scope}
The comparison is restricted to the Niemeier--Monster slice of the
Schellekens $c=24$ census.  The twenty-four Niemeier lattice VOAs are
class~$\mathbf G$ at the scalar shadow level, with
$\kappa_{\mathrm{latt}}=24$ and $r_{\max}=2$.  The unique $V_1=0$ row
represented by $\Vnat$ is class~$\mathbf M$, with
$\kapch(\Vnat)=12$ and $r_{\max}=\infty$.  The remaining forty-six
Schellekens rows require their own current algebra, scalar shadow, and
denominator data; the Niemeier and Monster rows do not determine them.
\end{remark}

\begin{remark}[Five-stratum archetype tagging of the moonshine input package]
\label{rem:moonshine-five-stratum-archetype}
\index{five-archetype dichotomy!moonshine input package}
\index{$\kappa$-stratification matrix!moonshine row}
This is an Open-quadrant statement for the $\Vnat,\Vsnat,W^{\Th}$
chart algebras, with the CY-coupled $\VKthree$ row present only via
the Vol.~III two-stage factorisation
$\Phi_2^{(\Sigma_1,C)}=\SpCh_{\Sigma_1,C}\circ\PhiFA_2$;
the presentation is chart-level, the Beilinson levels are $1$ and
$5$, and the hypotheses are the working loci of
Definition~\ref{def:moonshine-four-families} together with the
channel-supremum classification rule.

On the standard Open-quadrant Beilinson tower, the four chiral inputs
of Definition~\ref{def:moonshine-four-families} carry the following
archetype tags and rows in the
$5\times5$ $\kappa$-stratification matrix
$(\kapcat,\kappa^{\mathrm{Hodge}}_{\mathrm{ch}},
\kappa^{\mathrm{Heis}}_{\mathrm{ch}},\kapBKM,\kapfiber)$
with five distinct $\kappa$-measurements in each row. The
$\mathbf B$-row entry includes its pair $(\Sigma_{d-1},C)$ and is
not a one-stage $\Phi_d$.
\begin{center}
\renewcommand{\arraystretch}{1.3}
\begin{tabular}{l|c|p{0.45\linewidth}}
\toprule
\textbf{Object} & \textbf{Archetype} & \textbf{Five-stratum row} \\
\midrule
$V_{\Lambda_{24}}$ (Niemeier lattice VOA) & $\mathbf G$ &
 $(-,\,0,\,24,\,-,\,24)$: Heisenberg slot is the Leech rank;
 $r_{\max}=2$ \\
$\Vnat$ (FLM Monster module) & $\mathbf M$ &
 $(-,\,12,\,0,\,-,\,-)$: Hodge slot is $c/2{=}12$;
 $V_1{=}0$ kills the Heisenberg slot;
 $r_{\max}=\infty$ \\
$\Vsnat$ (Conway $N{=}1$ super-VOA) & $\mathbf M$ &
 $(-,\,6,\,0,\,-,\,-)$: Hodge slot is $c/2{=}6$;
 $r_{\max}=\infty$ \\
$W^{\Th}$ (Thompson graded module) & $\mathbf M$ (conditional) &
 row determined by an external chiral realisation $\VTh$ if supplied
 (Theorem~\ref{thm:thompson-chiral}) \\
$\VKthree$ ($c{=}6$ $\cN{=}4$ chiral sector) & $\mathbf B$
 (after Mukai enhancement) &
 anchored at $(0,\,0,\,3,\,5,\,24)$ on the
 Mukai-enhanced K3$\times E$ row \\
\bottomrule
\end{tabular}
\end{center}
The five entries are five distinct invariants on five distinct lanes.
A dash records that no Borcherds product, no compact CY fibre, or no
relevant slot is attached at the chart level. The Niemeier
$\mathbf G$ row and the FLM $\mathbf M$ row are different rows of the
matrix, indexed by the lattice rank (Heisenberg slot) and the Virasoro
$c/2$ scalar (Hodge slot), respectively.  The
class~$\mathbf B$ entry is reserved for the Mukai-enhanced K3$\times E$
witness whose canonical row is the class~$\mathbf B$ CY-host row;
$\VKthree$ alone is the chart algebra of an entry in that row only
after the Vol~III two-stage factorisation
$\Phi_2^{(\Sigma_1,C)}=\SpCh_{\Sigma_1,C}\circ\PhiFA_2$ has been
applied with $\Sigma_1$ a transverse one-real-dimensional manifold and
$C$ a fixed reference curve.
\end{remark}

\begin{remark}[Super sign convention]
\label{rem:moonshine-bar-super-conv}
For $\Vsnat$, $W^{\Th}$, and $\VKthree$ every bar primitive and every
trace is $\bZ/2$-graded.  The coefficient of weight $n$ is the
superdimension
\[
 a(n)=\sdim P_n=\dim P_n^{\bar 0}-\dim P_n^{\bar 1}.
\]
Odd collision OPEs remain collision OPEs; they are not replaced by
Laplace kernels.
\end{remark}

%%%%%%%%%%%%%%%%%%%%%%%%%%%%%%%%%%%%%%%%%%%%%%%%%%%%%%%%%%%%%%%%%%%%%%%%%%%%
\section{Primitive bar-Euler products}
\label{sec:moonshine-master}
%%%%%%%%%%%%%%%%%%%%%%%%%%%%%%%%%%%%%%%%%%%%%%%%%%%%%%%%%%%%%%%%%%%%%%%%%%%%

\begin{proposition}[Primitive bar-Euler product; \ClaimStatusProvedHere]
\label{prop:bar-euler-hilbert}
Let $A$ be an augmented weight-graded chiral algebra with finite
weight pieces, and let
\[
 P_A
 \;:=\;
 H^\bullet\!\bigl(\operatorname{Prim}\Barch(A)\bigr)
\]
be the primitive stratum of its chiral bar complex after taking the
weight-preserving bar differential.  Put
$a_A(n)=\sdim (P_A)_n$.  The primitive bar-Euler product is
\begin{equation}
\label{eq:bar-euler-hilbert}
 E_{\Barch}(A;q)
 \;:=\;
 \PE\!\left(\sum_{n\geq 1} a_A(n)q^n\right)
 \;=\;
 \prod_{n\geq 1}(1-q^n)^{-a_A(n)} .
\end{equation}
For a purely even object whose primitive bar homology is $\bar A$,
one has $a_A(n)=\dim \bar A_n$.
\end{proposition}

\begin{proof}
Work in the completed ring of formal power series in $q$.  Since each
weight piece $P_n$ is finite-dimensional, every coefficient of the
Euler product is a finite sum.  For a graded super vector space
$P=\bigoplus_{n\geq 1}P_n$ with finite weight pieces, the plethystic
exponential is defined by
\[
 \PE\!\left(\sum_{n\geq 1} a(n)q^n\right)
 =
 \exp\!\left(
   \sum_{r\geq 1}\frac{1}{r}\sum_{n\geq 1}a(n)q^{rn}
 \right).
\]
Taking the logarithm of the product on the right of
\eqref{eq:bar-euler-hilbert} gives the same series:
\[
 -\sum_{n\geq 1}a(n)\log(1-q^n)
 =
 \sum_{n,r\geq 1}\frac{a(n)}{r}q^{rn}.
\]
This is the symmetric, plethystic Euler product generated by the
primitive bar classes.

It is not the ordinary Euler characteristic of the tensor coalgebra
$T^c(s^{-1}\bar A)$.  Before taking a differential, the tensor
coalgebra has tensor-length series
$1/(1-\chi_q(s^{-1}\bar A))$ up to the suspension sign.  The product
\eqref{eq:bar-euler-hilbert} is obtained only after passing to the
primitive strata and then forming the free supercommutative
plethystic product.  Theorem~H is used exactly at this point in
Vol~I applications: finite cohomological amplitude allows the
primitive supertrace to be computed on homology without changing the
Euler characteristic.
\end{proof}

\begin{theorem}[Denominator/bar-Euler comparison criterion;
\ClaimStatusProvedHere]
\label{thm:moonshine-bar-euler-master}
\index{chiral moonshine!denominator comparison}
\index{Borcherds denominator identity!primitive bar-Euler product}
This is an Open-quadrant chain-level statement from Beilinson level
$2$ to level~$5$. The hypotheses are an augmented weight-graded
chiral algebra with finite weight pieces, a primitive-bar finite-type
model, and an external denominator-recognition datum as in
Definition~\ref{def:moonshine-finite-denominator-recognition}.

Let $A$ be an augmented weight-graded chiral algebra chosen in a
moonshine sector, and suppose an external generalized Kac--Moody
denominator is given in the form
\begin{equation}
\label{eq:moonshine-master}
 \mathcal D_A(\mathbf z)
 =
 e^\rho
 \prod_{\alpha\in\Delta_+}
 \bigl(1-e^{-\alpha}\bigr)^{\mu_A(\alpha)} .
\end{equation}
Let $\lambda:\Delta_+\to\bZ_{\geq 1}$ be a specialization of positive
roots to one weight variable, and let
$\varepsilon_A\in\{\pm1\}$ be the denominator-sign convention.  If the
specialized root multiplicities satisfy
\begin{equation}
\label{eq:moonshine-exponent-match}
 \sum_{\lambda(\alpha)=n}\mu_A(\alpha)
 =
 \varepsilon_A a_A(n)
 \qquad(n\geq 1),
\end{equation}
with $a_A(n)=\sdim(P_A)_n$ as in
Proposition~\ref{prop:bar-euler-hilbert}, then the normalized
specialized denominator is the signed inverse power of the primitive
bar-Euler product:
\[
 e^{-\rho}\mathcal D_A\!\mid_\lambda(q)
 =
 E_{\Barch}(A;q)^{-\varepsilon_A}.
\]
Thus denominator exponents equal to primitive supertraces give
$E_{\Barch}(A;q)^{-1}$, while reciprocal-product conventions give
$E_{\Barch}(A;q)$.
\end{theorem}

\begin{proof}
After specialization, \eqref{eq:moonshine-master} becomes
\[
 e^{-\rho}\mathcal D_A\!\mid_\lambda(q)
 =
 \prod_{n\geq 1}
 (1-q^n)^{\sum_{\lambda(\alpha)=n}\mu_A(\alpha)} .
\]
Substituting \eqref{eq:moonshine-exponent-match} gives
\[
 \prod_{n\geq 1}(1-q^n)^{\varepsilon_A a_A(n)}
 =
 E_{\Barch}(A;q)^{-\varepsilon_A}
\]
by Proposition~\ref{prop:bar-euler-hilbert}.  No chiral algebra
produces a Borcherds denominator from this formal calculation; the
denominator, sign convention, root specialization, and exponent match
are external inputs.
\end{proof}

\begin{definition}[Finite denominator-recognition datum]
\label{def:moonshine-finite-denominator-recognition}
For a sector $(A,G_A)$ and an external denominator
\eqref{eq:moonshine-master}, a finite denominator-recognition datum on
a downward-saturated finite root window $W\subset\Delta_+$ consists of:
\begin{enumerate}[label=\textup{(\roman*)}]
\item a finite-dimensional, $G_A$-stable chain model for
$H^\bullet(\operatorname{Prim}\Barch(A))$ in the weights represented by
$W$;
\item the root specialization $\lambda$, Weyl-vector normalization,
parity convention, and sign $\varepsilon_A$;
\item a parity-preserving comparison map from this primitive bar model
to the corresponding finite sum of BKM root spaces such that
\[
 \sum_{\substack{\alpha\in W\\ \lambda(\alpha)=n}}\mu_A(\alpha)
 =
 \varepsilon_A\,
 \Str\!\left(g\mid H^\bullet(\operatorname{Prim}\Barch(A))_n\right)
\]
for every represented weight $n$ and every twining element for which a
twined comparison is part of the datum;
\item if the datum includes a source BKM algebra, rather than only a
scalar target product, the finite-window Chevalley, real-Serre, super-Jacobi,
radical, and PBW compatibility statements for the same window.
\end{enumerate}
Passing from the finite windows to a completed denominator identity
requires the corresponding Mittag--Leffler or finite-determination
bound.  Without these data, an Igusa or Borcherds product is only an
external automorphic target.
\end{definition}

\begin{remark}[Finite type, Verdier duals, and centres]
\label{rem:moonshine-finite-type-verdier-firewall}
The Euler products use finite-dimensional primitive bar homology in
each weight.  A Verdier dual is defined only under the
finite-type, perfect, or continuous-dual hypotheses of the cited
bar--cobar theorem.  The objects
\[
 A,\qquad
 \Barch(A),\qquad
 A^{\mathrm i}=H^\bullet\Barch(A),\qquad
 A^!,\qquad
 \ZcSig(A)
\]
remain distinct: $\Omega\Barch(A)\simeq A$ is bar--cobar inversion;
$A^!$ is obtained only after Verdier or continuous linear duality from
$A^{\mathrm i}$; and
$\ZcSig(A)=C^\bullet_{\mathrm{ch}}(A,A)$ is the Hochschild/bulk
object, not a bar coalgebra.
\end{remark}

\begin{remark}[Sectorwise exponent data]
\label{rem:moonshine-sectorwise-exponent-data}
The comparison criterion has four different inputs.
\begin{center}
\begin{tabular}{c|p{0.34\linewidth}|p{0.37\linewidth}}
\textbf{sector} & \textbf{external product or trace} & \textbf{bar comparison}\\
\hline
$\Vnat$ &
Borcherds' monster Lie algebra denominator
$J(\sigma)-J(\tau)
=p^{-1}\prod_{m>0,n\in\bZ}(1-p^m q^n)^{c(mn)}$ &
The exponents are coefficients of
$J(q)=q^{-1}+\sum_{n\geq 1}c(n)q^n$ and
$\dim \Vnat_N=c(N-1)$.  The shift and the two variables are part of
the theorem; setting $p=0$ is not a specialization of the product.\\
$\Vsnat$ &
Duncan's $\Co_0$-twined supertraces on $V^{s\natural}$ &
These traces determine primitive superdimensions for a Conway
bar-Euler product.  A Conway BKM denominator may be compared only
after its root multiplicities are identified with those
superdimensions.\\
$W^{\Th}$ &
Harvey--Rayhaun / Griffin--Mertens Thompson graded traces &
The literature supplies a graded $\Th$-module and modular functions.
A chiral algebra $\VTh$ and a denominator-root match are additional
data, not consequences of the Conway construction.\\
$\VKthree$ &
The K3 elliptic genus and the DMVV--Gritsenko--Nikulin multiplicative
lift to the Igusa product &
The exponents are Fourier--Jacobi coefficients
$c(4mn-\ell^2)$ of the K3 elliptic genus.  They belong to the
second-quantized K3 product, not to the bare Hilbert series of the
$c=6$ $\cN{=}4$ vacuum algebra.  A bar comparison requires the finite
denominator-recognition datum of
Definition~\ref{def:moonshine-finite-denominator-recognition}.
\end{tabular}
\end{center}
\end{remark}

\begin{remark}[Four inputs, not one functor]
\label{rem:moonshine-four-inputs-not-one-functor}
The comparison above uses four distinct sources.  $\Vnat$ is the FLM
Leech $\bZ/2$ orbifold.  $\Vsnat$ is Duncan's Conway super-VOA.
$W^{\Th}$ is a Thompson graded module unless a separate chiral model
is supplied.  $\VKthree$ comes from the K3 $\cN{=}4$ elliptic-genus
sector.  The two-stage CY-to-chiral factorization
$\Phi_2=\SpCh_{\Sigma_1,C}\circ\PhiFA_2$ applies to the K3 entry; it
does not manufacture the Monster, Conway, and Thompson entries.
\end{remark}

%%%%%%%%%%%%%%%%%%%%%%%%%%%%%%%%%%%%%%%%%%%%%%%%%%%%%%%%%%%%%%%%%%%%%%%%%%%%
\section{The Monster sector}
\label{sec:vnat-e3-topological}
%%%%%%%%%%%%%%%%%%%%%%%%%%%%%%%%%%%%%%%%%%%%%%%%%%%%%%%%%%%%%%%%%%%%%%%%%%%%

\begin{definition}[E$_3$-topological descent through the Leech orbifold]
\label{def:e3-topological}
For the Monster sector, an E$_3$-topological statement means the
finite-orbifold descent statement from the Leech lattice
holomorphic-topological model to the FLM boundary VOA.  The data are:
\begin{enumerate}[label=\textup{(\roman*)}]
\item the abelian Leech lattice model with its chain-level
E$_3$-topological structure;
\item a lifted $\bZ/2$ action with trivial Dijkgraaf--Witten anomaly
class in $H^3(B\bZ/2;U(1))$;
\item identification of the descended boundary chiral algebra with
$V_\Lambda^+\oplus V_\Lambda^{\mathrm{tw},+}=\Vnat$.
\end{enumerate}
This is a statement about the descended HT model.  It is not an
identification of $A=\Vnat$, $\Barch(A)$,
$A^{\mathrm i}=H^\bullet\Barch(A)$, the Verdier Koszul dual $A^!$, or
the derived chiral centre $\ZcSig(A)$.  The equivalence
$\Omega\Barch(A)\simeq A$ is bar--cobar inversion, not Verdier
Koszul duality and not a bulk computation.
\end{definition}

\begin{theorem}[$\Vnat$ from the Leech orbifold E$_3$ descent;
\ClaimStatusConditional]
\label{thm:v-natural-e3-topological}
\index{Monster module!E_3-topological descent}
\index{Leech orbifold!BV anomaly}
This is an Open-quadrant chain-level statement from Beilinson level
$0$ to level~$1$ through finite-orbifold descent. Its hypotheses are
the three data of Definition~\ref{def:e3-topological}: the abelian
Leech HT model, an anomaly-free $\bZ/2$ lift, and the FLM boundary
identification. The class transition is recorded separately in
Remark~\ref{rem:vnat-class-transition-clarified}.

Assume the three data of Definition~\ref{def:e3-topological}.  Then
the descended finite-orbifold HT factorisation algebra has a
chain-level E$_3$-topological structure and boundary chiral algebra
$\Vnat$.  The conformal vector of $\Vnat$ is the FLM conformal vector
of central charge $24$ inherited through the Leech orbifold; since
$\Vnat_1=0$, it is not a Sugawara tensor built from weight-$1$
currents.
\end{theorem}

\begin{proof}
FLM construct
$\Vnat=V_\Lambda^+\oplus V_\Lambda^{\mathrm{tw},+}$ and prove the
twisted Jacobi identity for this orbifold VOA.  The boundary
identification in Definition~\ref{def:e3-topological}(iii) imports
that VOA as the boundary of the descended HT theory.  The abelian
Leech E$_3$ operations descend through finite orbifolding precisely
when the BV associator anomaly is trivial; this is
Definition~\ref{def:e3-topological}(i)--(ii).  Therefore the
descended HT factorisation algebra is E$_3$-topological, and its
boundary is $\Vnat$.

The conformal-vector clause is FLM's construction of the
Moonshine module.  The weight-two space decomposes as the conformal
line together with the Griess algebra, but the absence of
weight-one currents forbids a Heisenberg or affine Sugawara formula
inside $\Vnat$.
\end{proof}

\begin{remark}[Monster denominator and scalar bar data]
\label{rem:vnat-bar-euler-status}
Borcherds' denominator for the monster Lie algebra has exponents
$c(mn)$ from $J(q)=q^{-1}+\sum c(n)q^n$.  The coefficient shift
$\dim \Vnat_N=c(N-1)$ is essential.  The denominator formula, the
FLM construction, and the primitive bar-Euler product meet only after
this shift and the chosen root specialization are recorded.  The
scalar shadow data of Proposition~\ref{prop:moonshine-kappa} are
\[
 \kapch(\Vnat)=12,\qquad
 S_4^{\mathrm{Vir}}=\frac{5}{1704},\qquad
 \Delta=\frac{20}{71},\qquad
 r^{\mathrm{Vir}}_{24}(z)=\frac{12}{z^3}+\frac{2T}{z}.
\]
These data distinguish $\Vnat$ from the Leech lattice VOA at the
scalar level, but they do not compute the Monster action on
$H^\bullet(\operatorname{Prim}\Barch(\Vnat))$.  A Monster-equivariant
bar denominator statement requires the finite recognition datum of
Definition~\ref{def:moonshine-finite-denominator-recognition}, with
graded Monster traces on primitive bar homology.  The E$_3$ descent
theorem is independent of that denominator calculation.
\end{remark}

\begin{remark}[Shadow class versus E$_3$ descent]
\label{rem:vnat-class-transition-clarified}
Chapter~\ref{chap:moonshine}
Remark~\ref{rem:moonshine-orbifold-class-transition} records that
the Leech orbifold changes the shadow class from $\mathbf G$ to
$\mathbf M$.  E$_3$ descent is a different invariant: it is transported
from the abelian Leech HT model through an anomaly-free finite
orbifold.  Thus the class transition does not by itself prove or
disprove E$_3$-topological descent.
\end{remark}

%%%%%%%%%%%%%%%%%%%%%%%%%%%%%%%%%%%%%%%%%%%%%%%%%%%%%%%%%%%%%%%%%%%%%%%%%%%%
\section{The Conway and Thompson sectors}
\label{sec:conway-thompson}
%%%%%%%%%%%%%%%%%%%%%%%%%%%%%%%%%%%%%%%%%%%%%%%%%%%%%%%%%%%%%%%%%%%%%%%%%%%%

\begin{theorem}[Conway chiral input; \ClaimStatusProvedElsewhere]
\label{thm:conway-chiral-structure}
\index{Conway module!chiral structure}
\index{Duncan module!primitive bar-Euler product}
This is an Open-quadrant chain-level statement at Beilinson levels
$1$ and~$2$. The hypotheses are Duncan's $N{=}1$ super-VOA structure
on $\Vsnat$ and the $\Co_0$ action by super-VOA automorphisms; no
Conway BKM denominator is asserted.

Duncan's Conway module $\Vsnat$ is a $\bZ/2$-graded
$N{=}1$ super-VOA of central charge $12$ with
$\dim \Vsnat_{1/2}=0$ and a $\Co_0$ action by super-VOA
automorphisms.  Its primitive bar-Euler product is
\[
 E_{\Barch}(\Vsnat;q)
 =
 \prod_{n\geq 1}(1-q^n)^{-a_{\Vsnat}(n)},
 \qquad
 a_{\Vsnat}(n)=\sdim H^\bullet(\operatorname{Prim}\Barch(\Vsnat))_n .
\]
For $g\in\Co_0$, replacing $a_{\Vsnat}(n)$ by the primitive
supertrace $\Str(g\mid P_{\Vsnat,n})$ gives the twined bar-Euler
product.  This theorem does not assert a Conway BKM denominator
unless an external denominator with the same exponents is supplied.
\end{theorem}

\begin{proof}
Duncan's construction supplies the $N{=}1$ super-VOA, the central
charge $12$, the vanishing of the degree-$1/2$ space, and the
$\Co_0$ action.  Since $\Co_0$ preserves the super-VOA structure, it
preserves the bar differential and acts on the primitive bar
homology.  Proposition~\ref{prop:bar-euler-hilbert} gives the
untwined and twined products by taking superdimensions or
supertraces on those primitive weight spaces.
\end{proof}

\begin{theorem}[Thompson chiral realization criterion;
\ClaimStatusConditional]
\label{thm:thompson-chiral}
\index{Thompson moonshine!chiral realization criterion}
This is an Open-quadrant chain-level statement at Beilinson levels
$1$ and~$2$. It assumes an external $\Th$-equivariant chiral
superalgebra $\VTh$ together with an isomorphism of graded
$\Th$-modules
$H^\bullet(\operatorname{Prim}\Barch(\VTh))_n\cong W^{\Th}_n$;
a Thompson BKM denominator enters only if its specialised root
multiplicities match the named graded supertraces.

Let $W^{\Th}=\bigoplus_n W^{\Th}_n$ be the graded Thompson module
whose traces are the Harvey--Rayhaun Thompson moonshine functions.
Suppose a $\Th$-equivariant chiral superalgebra $\VTh$ is given
together with an isomorphism of graded $\Th$-modules
\[
 H^\bullet(\operatorname{Prim}\Barch(\VTh))_n
 \;\cong\;
 W^{\Th}_n
\]
for all $n$, compatible with parity.  Then
\[
 E_{\Barch}^{g}(\VTh;q)
 =
 \prod_{n\geq 1}
 (1-q^n)^{-\Str(g\mid W^{\Th}_n)}
 \qquad(g\in\Th).
\]
If, in addition, a Thompson BKM denominator has specialized root
multiplicities $-\Str(g\mid W^{\Th}_n)$, the denominator/bar
comparison of Theorem~\ref{thm:moonshine-bar-euler-master} applies.
\end{theorem}

\begin{proof}
The first displayed identity is exactly
Proposition~\ref{prop:bar-euler-hilbert} after the assumed
identification of primitive bar homology with the Thompson module.
The denominator clause is Theorem~\ref{thm:moonshine-bar-euler-master}.
\end{proof}

\begin{remark}[No fixed-point construction asserted]
\label{rem:thompson-as-conway-subsector}
The data above are weaker than a construction of $\VTh$ as a
$3A$-fixed-point subalgebra of $\Vsnat$.  The Thompson group is not
obtained by taking Conway invariants, and no rank-halving or
$\kapch(\VTh)$ computation is made.  Harvey--Rayhaun and
Griffin--Mertens supply graded Thompson moonshine functions; a
compatible chiral algebra model is the extra input named in
Theorem~\ref{thm:thompson-chiral}.
\end{remark}

%%%%%%%%%%%%%%%%%%%%%%%%%%%%%%%%%%%%%%%%%%%%%%%%%%%%%%%%%%%%%%%%%%%%%%%%%%%%
\section{The Mathieu sector and the K3/Mukai separation}
\label{sec:mathieu-k3-link}
%%%%%%%%%%%%%%%%%%%%%%%%%%%%%%%%%%%%%%%%%%%%%%%%%%%%%%%%%%%%%%%%%%%%%%%%%%%%

\begin{theorem}[Mathieu elliptic genus and second quantization]
\label{thm:mathieu-class-from-N4-K3}
\ClaimStatusConditional
\index{Mathieu moonshine!K3 elliptic genus}
\index{Mathieu moonshine!Borcherds lift}
This is an Open-quadrant assertion for the chart algebra $\VKthree$,
with a CY-quadrant comparison only after the Vol.~III two-stage
factorisation
$\Phi_2^{(\Sigma_1,C)}=\SpCh_{\Sigma_1,C}\circ\PhiFA_2$ with
$\Sigma_1$ a transverse one-real-dimensional manifold and $C$ a
fixed reference curve. It is stated in the chain presentation at
Beilinson levels $1$ and $5$. Its hypotheses are the
Eguchi--Ooguri--Tachikawa decomposition into $c=6$ $\cN{=}4$
characters, Gannon's $\Mtw$-module theorem, the DMVV multiplicative
lift, and the Gritsenko--Nikulin Borcherds lift; for the bar
comparison one also assumes the finite
denominator-recognition datum of
Definition~\ref{def:moonshine-finite-denominator-recognition}.

Let $Z_{K3}(\tau,z)$ be the K3 elliptic genus, decomposed into
$c=6$ $\cN{=}4$ characters as in Eguchi--Ooguri--Tachikawa and
Gannon.  Then:
\begin{enumerate}[label=\textup{(\roman*)}]
\item Gannon's theorem gives a graded $\Mtw$-module whose traces are
the Mathieu twining functions.  This is a module theorem; it is not
a full $\Mtw$ action on one K3 sigma model.
\item The DMVV multiplicative lift and the
Gritsenko--Nikulin denominator identify the second-quantized K3
elliptic genus with an Igusa/Siegel automorphic product.  Its
exponents are Fourier--Jacobi coefficients
$c(4mn-\ell^2)$ of $Z_{K3}$.
\item The primitive bar-Euler product of the bare algebra
$\VKthree$ uses the primitive superdimensions of
$\Barch(\VKthree)$.  It is compared with the DMVV/Gritsenko--Nikulin
product only after a chiral model has been fixed whose primitive
supertrace is the same Fourier--Jacobi coefficient system, together
with the sign, Weyl-vector, parity, and finite-window recognition data
of Definition~\ref{def:moonshine-finite-denominator-recognition}.
\end{enumerate}
\end{theorem}

\begin{proof}
Clause (i) is Gannon's proof of Mathieu moonshine.  Clause (ii) is
the DMVV second-quantized elliptic-genus product together with the
Gritsenko--Nikulin Borcherds lift.  The product side counts
Fourier--Jacobi coefficients of the elliptic genus; those are
second-quantized K3 coefficients, not weight dimensions of the
abstract small $\cN{=}4$ vacuum algebra.  Clause (iii) is therefore
the signed coefficient-matching and finite-recognition condition of
Theorem~\ref{thm:moonshine-bar-euler-master} applied to the K3
sector.
\end{proof}

\begin{remark}[K3 values are not interchangeable]
\label{rem:mathieu-k3-values-separated}
The identities
\[
 \kapch(\VKthree)=3,\qquad
 \kapcat(K3)=2,\qquad
 \kapBKM(\mathfrak g_{\Delta_5})=5,\qquad
 K^\kappa(\mathcal H_{\mathrm{Muk}}(K3))=8,\qquad
 \kapfiber(\widetilde H(K3,\bZ))=24
\]
belong to five constructions: Virasoro curvature, CY$_2$ categorical
Euler characteristic, Borcherds product weight, Mukai-enhanced scalar
conductor, and Mukai-lattice rank.  The Igusa product uses the
Borcherds/DMVV surface.  The celestial K3 construction uses the
$\cN{=}4$ matter surface.  The scalar conductor $K^\kappa=8$ is the
Mukai-doubling complementarity value of
Proposition~\ref{prop:archetype-complementarity-bridge}, not the BKM
weight~$5$.  A bar-Euler identity between these surfaces requires an
explicit finite recognition datum; it is not a formal consequence of
sharing the word ``K3''.
\end{remark}

%%%%%%%%%%%%%%%%%%%%%%%%%%%%%%%%%%%%%%%%%%%%%%%%%%%%%%%%%%%%%%%%%%%%%%%%%%%%
\section{Twining, shadow towers, and congruence tests}
\label{sec:shadow-tower-twining}
%%%%%%%%%%%%%%%%%%%%%%%%%%%%%%%%%%%%%%%%%%%%%%%%%%%%%%%%%%%%%%%%%%%%%%%%%%%%

\begin{theorem}[Twined shadow comparison: K3 domain and general
criterion; \ClaimStatusConditional]
\label{thm:shadow-tower-twining-universality}
\index{shadow tower!twining comparison}
\index{McKay--Thompson series!primitive bar traces}
This is an Open-quadrant chain-level comparison from Beilinson level
$2$ to the level-$5$ scalar shadow. For K3 it uses the Vol.~II
celestial--Mathieu theorem and the Mellin--shadow dictionary. For a
general pair $(A,G_A)$ it assumes a chain map compatible with
twining and depth filtrations from the primitive bar complex to the
Rademacher polar-data complex, together with the sign, Weyl-vector,
parity, and finite coefficient bounds of
Definition~\ref{def:moonshine-finite-denominator-recognition}.

For the K3 celestial algebra, Vol~II proves a Mathieu statement:
the $g$-twined shadow-tower coefficients of the K3 matter factor
match the Rademacher coefficients of the corresponding Mathieu mock
modular form $h_g(\tau)$.  For a general moonshine input
$(A,G_A)$, the analogous identity
\begin{equation}
\label{eq:shadow-tower-mck-thompson}
S_r(A)^g
\;=\;
\mathrm{Rad}_r\!\bigl[T^{G_A}_g\bigr]
\end{equation}
holds only under an additional comparison map from the depth-$r$
bar-cohomology primitive to the polar/Rademacher datum of the
twined modular function $T^{G_A}_g$.
\end{theorem}

\begin{proof}
The K3 clause is the Vol~II celestial--Mathieu theorem together
with the Mellin--shadow dictionary: the celestial K3 matter factor
has twined mock modular forms $h_g$, and the soft/Mellin residue
identifies the relevant shadow coefficient with the Rademacher
coefficient.  This proof uses the tensor decomposition of the
celestial K3 algebra and the Gannon $\Mtw$-module.  The $\Mtw$ action
is an action on the Mathieu coefficient module; it is not a geometric
$\Mtw$ action on the $24$ marked points of a class-$\mathcal S$
sphere.  A passage through the $\Sigma_{0,24}$ AGT or Schur-index
realisation uses only the finite Schur--Igusa recognition data; a
cyclotomic marking gives $\mu_{24}$ symmetry, not Mathieu symmetry.

For Monster, Conway, and Thompson, equation
\eqref{eq:shadow-tower-mck-thompson} has two independent sides:
the left side is a bar-cohomology invariant of a specified chiral
object, while the right side is analytic data of a Hauptmodul or
mock modular form.  Linearity of traces is not enough to identify
them.  The missing datum is the finite recognition package: a chain
map compatible with twining and depth filtrations from the primitive
bar complex to the Rademacher polar-data complex, together with the
sign, Weyl-vector, parity, and finite coefficient bounds required by
Definition~\ref{def:moonshine-finite-denominator-recognition}.  Once
such data are supplied, the equality is a theorem by functoriality of
the Mellin/Rademacher construction; before that, it is a criterion.
\end{proof}

\begin{corollary}[Kummer congruence test, not automatic congruence;
\ClaimStatusConditional]
\label{cor:kummer-congruence-moonshine}
\index{Kummer congruence!moonshine test}
Fix a sector $(A,G_A)$ and $g\in G_A$.  Suppose the twined series
$T^{G_A}_g$ admits a decomposition
\[
 T^{G_A}_g = E_g + C_g
\]
in which $E_g$ is an Eisenstein series with integral $q$-coefficients
and $C_g$ is a cusp or mock part whose relevant coefficients vanish
modulo a prime $p$.  If $p=691$ in weight $12$ or $p=3617$ in weight
$16$, then the Eisenstein part satisfies the usual Kummer congruence
modulo $p$, and the same congruence holds for the selected coefficient
of $T^{G_A}_g$.
\end{corollary}

\begin{proof}
The primes $691$ and $3617$ divide the numerators of the Bernoulli
numbers $B_{12}$ and $B_{16}$, respectively.  Kummer congruences for
Eisenstein series give the stated congruence on $E_g$.  The
hypothesis on $C_g$ is exactly what is needed to transfer the
congruence from $E_g$ to the full twined series.  Without this
hypothesis no congruence for a McKay--Thompson or Mathieu twining
series follows from the bar formalism alone.
\end{proof}

\begin{remark}[Finite computation required]
\label{rem:moonshine-falsification-test}
For a concrete sector and conjugacy class the test is finite:
compute the relevant twined coefficients from the moonshine character
table, compute the bar primitive trace from the chosen chiral model,
and compare them modulo $691$ or $3617$ after the Eisenstein/cusp
decomposition has been fixed.  The statement is a verification
problem for the chosen sector, not an automatic consequence of the
bar formalism.
\end{remark}

%%%%%%%%%%%%%%%%%%%%%%%%%%%%%%%%%%%%%%%%%%%%%%%%%%%%%%%%%%%%%%%%%%%%%%%%%%%%
\section*{Signed denominator bridge}
%%%%%%%%%%%%%%%%%%%%%%%%%%%%%%%%%%%%%%%%%%%%%%%%%%%%%%%%%%%%%%%%%%%%%%%%%%%%

\noindent
The signed exponent match is the invariant.  Borcherds
products, DMVV products, Conway supertraces, Thompson graded traces,
and chiral bar primitives are different objects.  When a cited
denominator product has exponents equal to
$\varepsilon_A$ times the primitive graded supertrace of the chosen
chiral object, Theorem~\ref{thm:moonshine-bar-euler-master} identifies
the normalized product with
$E_{\Barch}(A;q)^{-\varepsilon_A}$.  The Monster and K3 sectors have
named product formulae with known coefficient systems; the Conway and
Thompson sectors require the extra denominator or chiral-realization
data explicitly stated above.

%=======================================================================
% End of chiral_moonshine_unified.tex
%=======================================================================
